% from ega4-10-fr.tex
\oldpage[IV-3]{98}
Supposons d'abord que $f(X_0)$ ne soit pas très dense dans $X$~; il existe
alors deux parties fermées distinctes $Y\supset Y'$ de $X$ telles que
$Y\cap f(X_0)=Y'\cap f(X_0)$. Soit $\mathcal F$ (resp. $\mathcal F'$) le
faisceau de groupes abéliens sur $X$ image directe par l'injection canonique
$Y\to X$ (resp. $Y'\to X$) du faisceau simple associé au préfaisceau constant
$\mathbf Z$ sur $Y$ (resp. $Y'$)~; la définition de $f^*$
($\mathbf 0_{\mathrm I}$, 3.7.1) montre que $f^*(\mathcal F)$ est isomorphe à
$f^*(\mathcal F')$ mais $\mathcal F$ n'est pas isomorphe à $\mathcal F'$,
donc la condition \emph{c)} n'est pas vérifiée. (On notera que le foncteur
$f^*$ n'est même pas alors \emph{fidèle}, car si $u$ et $v$ sont
l'automorphisme identique et l'endomorphisme nul de $\mathcal F/\mathcal F'$,
$f^*(u)$ et $f^*(v)$ sont égaux.)

Montrons maintenant que \emph{c)} entraîne que la topologie de $X_0$ est
l'image réciproque de celle de $X$ par $f$. Notons que si la condition
\emph{c)} est vérifiée pour la catégorie des faisceaux d'ensembles, elle l'est
aussi pour la catégorie des faisceaux de groupes abéliens, car il résulte
aussitôt des définitions ($\mathbf 0_{\mathrm{III}}$, 8.2.5) que cette dernière
n'est autre que la catégorie des \emph{objets en groupes abéliens} dans la
catégorie des faisceaux d'ensembles. Il suffira donc de prouver notre assertion
en supposant \emph{c)} vérifiée pour les catégories de faisceaux de groupes
abéliens sur $X$ et $X_0$. Or, si $\mathbf Z_X$ désigne le faisceau simple sur
$X$ associé au préfaisceau constant égal à $\mathbf Z$, on a un isomorphisme
canonique
\[
\Gamma(X,\mathcal F)\xrightarrow{\sim}\operatorname{Hom}(\mathbf Z_X,\mathcal F)
\]
fonctoriel en $\mathcal F$ (même raisonnement que dans
($\mathbf 0_{\mathrm I}$, 5.1.1))~; comme il est clair par définition
($\mathbf 0_{\mathrm I}$, 3.7.1) que $f^*(\mathbf Z_X)=\mathbf Z_{X_0}$,
l'hypothèse que $f^*$ est une équivalence entraîne que l'homomorphisme
canonique $\Gamma(X,\mathcal F)\to\Gamma(X_0,f^*(\mathcal F))$ est
\emph{bijectif} pour tout faisceau de groupes abéliens $\mathcal F$ sur $X$.
Or, soit $Y_0$ une partie fermée de $X_0$~; soit $\mathcal F_0$ le faisceau
sur $X$ image directe par l'injection canonique $Y_0\to X_0$ du faisceau
simple associé au préfaisceau constant $\mathbf Z$ sur $Y_0$. Comme $f^*$ est
une équivalence, $\mathcal F_0$ est isomorphe à un faisceau de la forme
$f^*(\mathcal F)$, où $\mathcal F$ est un faisceau de groupes abéliens sur
$X$. Considérons alors la section $s_0$ de $\mathcal F_0$ au-dessus de $X_0$
telle que $(s_0)_{x_0}=1_{x_0}$ si $x_0\in Y_0$,
$(s_0)_{x_0}=0_{x_0}$ si $x_0\notin Y_0$. Les remarques précédentes montrent
qu'il existe une section $s$ de $\mathcal F$ au-dessus de $X$ telle que
$(s_0)_{x_0}=s_{f(x_0)}$ pour tout $x_0\in X_0$ (les fibres
$(\mathcal F_0)_{x_0}$ et $\mathcal F_{f(x_0)}$ étant canoniquement
identifiées ($\mathbf 0_{\mathrm I}$, 3.7.1))~; cela entraîne que $Y_0$ est
l'image réciproque par $f$ de l'ensemble $Y$ des $x\in X$ tels que
$s_x\neq0_x$, et $y$ est fermé dans $X$. C.Q.F.D.

\begin{definition}[10.2.2]
\label{IV.10.2.2-fr}
Lorsqu'une application continue $f:X_0\to X$ vérifie les conditions
équivalentes de (10.2.1), on dit que $f$ est un \emph{quasi-homéomorphisme} de
$X_0$ dans $X$.
\end{definition}

En vertu de (10.2.1, \emph{b)}), dire qu'une partie $X_0$ d'un espace
topologique $X$ est \emph{très dense} dans $X$ signifie donc que l'injection
canonique $X_0\to X$ est un \emph{quasi-homéomorphisme}.

\begin{corollary}[10.2.3]
\label{IV.10.2.3-fr}
Le composé de deux quasi-homéomorphismes est un quasi-homéomorphisme.
\end{corollary}

Cela résulte aussitôt de (10.2.1, \emph{a)}).

\begin{corollary}[10.2.4]
\label{IV.10.2.4-fr}
Soient $f:X\to Y$ un quasi-homéomorphisme, $Y'$ une partie localement
quasi-constructible de $Y$, $X'=f^{-1}(Y')$~; alors la restriction
$f'=f|X':X'\to Y'$ est un quasi-homéomorphisme.
\end{corollary}

Il est clair en vertu de (10.2.1, \emph{b)}) que la topologie induite sur $X'$
par celle de $X$ est l'image réciproque par $f'$ de la topologie induite sur
$Y'$ par celle de $Y$. D'autre part,
